\subsection{Criteria for integrability}

THEOREM 5. --- In order that a mapping f of X into a Banach space F be p-th power integrable ( \(1 \leqslant p < +\infty\) ), it is necessary and sufficient that f be measurable and that \(N_{p}(\mathbf{f})\) be finite.

The condition is necessary: for, if \(f \in L_{F}^{p}\) then there exists a sequence \((\mathbf{g}_{n})\) of continuous functions with compact support that converges almost everywhere to f (§3, No.~4, Cor. 2 of Th. 3); by Th. 2 of No.~4, f is measurable.

To prove that the conditions are sufficient, we first establish a lemma:

Lemma 1. --- Let g be a function with values in F, such that \(\mathrm{N}_{p}(\mathbf{g}) < +\infty\) (in other words, a function in \(F_{F}^{p}\) ). The set A of points \(x \in X\) such that \(\mathbf{g}(x) \neq 0\) is contained in the union of a negligible set and a sequence of compact sets.

Let \(A_{n}\) be the set of points \(x \in X\) such that \(|\mathbf{g}(x)| \geqslant 1/n\) ; A is the union of the \(A_{n}\) , and \(\varphi_{A_{n}} \leqslant n|\mathbf{g}|\) , whence \(|\mu|^{*}(A_{n}) \leqslant (nN_{p}(\mathbf{g}))^{p}\) ; it follows that \(A_{n}\) is contained in the union of a negligible set and a sequence of compact sets (§4, No.~6, Cor. 3 of Th. 4), therefore so is A.

The lemma proved, consider first the case that \(\mathbf{f}\) has compact support \(\mathbf{K}\) . By Cor. 1 of Th. 3 of No.~5, there exists a sequence \((\mathbf{g}_n)\) of measurable step functions such that \(|\mathbf{g}_n(x)| \leqslant |\mathbf{f}(x)|\) at every point \(x \in \mathbf{X}\) and such that \(\mathbf{g}_n(x)\) tends to \(\mathbf{f}(x)\) almost everywhere. Now, \(\mathbf{g}_n\) is a linear combination of characteristic functions of measurable sets contained in \(\mathbf{K}\) ; since these sets are integrable by Prop. 3 of No.~1, \(\mathbf{g}_n\) belongs to \(\mathcal{L}_{\mathrm{F}}^p\) . Since \(\mathrm{N}_p(\mathbf{f}) < +\infty\) , Lebesgue's theorem (§3, No.~7, Th. 6) shows that \(\mathbf{f}\) belongs to \(\mathcal{L}_{\mathrm{F}}^p\) .

In the general case, it follows from Lemma 1 that there exists an increasing sequence \((\mathrm{K}_n)\) of compact sets such that \(\mathbf{f}(x)\) is zero almost everywhere in the complement of the union of the \(\mathrm{K}_n\) . Let \(\mathbf{f}_n\) be the function equal to \(\mathbf{f}(x)\) on \(\mathrm{K}_n\) and to 0 elsewhere; \(\mathbf{f}_n\) is measurable by No.~3, Cor. 5 of Th. 1; since \(|\mathbf{f}_n| \leqslant |\mathbf{f}|\) , \(\mathbf{f}_n\) belongs to \(\mathcal{L}_{\mathrm{F}}^p\) by the first part of the argument. Since \(\mathbf{f}(x)\) is equal almost everywhere to the limit of the sequence \(\mathbf{f}_n(x)\) , Lebesgue's theorem again proves that \(\mathbf{f} \in \mathcal{L}_{\mathrm{F}}^p\) , which completes the proof.

One should take care to note that a function that is locally negligible but not negligible is not integrable; thus, a function equal locally almost everywhere to an integrable function is not necessarily integrable.

COROLLARY 1. --- For a set to be integrable, it is necessary and sufficient that it be measurable and have finite outer measure.

COROLLARY 2. --- Let \((\mathrm{F}_{n})\) be a sequence of topological spaces; for every index n, let \(f_{n}\) be a measurable mapping of X into \(F_{n}\) and let \(f(x) = (f_{n}(x)) \in \mathrm{F} = \prod_{n} \mathrm{F}_{n}\) ; finally, let u be a continuous mapping of \(f(\mathrm{X})\) into a Banach space G. For the function \(\mathbf{g}(x) = \mathbf{u}(f(x))\) to be integrable, it is necessary and sufficient that \(\mathrm{N}_{1}(\mathbf{g}) < +\infty\) .

For, \(\mathbf{g}\) is measurable (No.~3, Th. 1).

COROLLARY 3. --- For every integrable function f and every measurable set A, the function \(f\varphi_{A}\) is integrable.

For, it follows from Th. 5 and No.~3, Cor. 5 of Th. 1 that \(\mathbf{f}\varphi_{\mathrm{A}}\) is measurable, and \(\mathrm{N}_1(\mathbf{f}\varphi_{\mathrm{A}}) \leqslant \mathrm{N}_1(\mathbf{f})\) .

We write \(\int_{\mathrm{A}} \mathbf{f} d\mu = \int \mathbf{f} \varphi_{\mathrm{A}} d\mu\) (or \(\int_{\mathrm{A}} \mathbf{f} \mu\) ) for every integrable function \(\mathbf{f}\) and every measurable set \(\mathrm{A}\) . We also write \(\int_{\mathrm{A}}^{*} f d|\mu|\) (or \(\int_{\mathrm{A}}^{*} f |\mu|\) ) in place of \(\int^{*} f \varphi_{\mathrm{A}} d|\mu|\) for every numerical function (finite or not) \(f \geqslant 0\) (on setting \(f(x) \varphi_{\mathrm{A}}(x) = 0\) if \(f(x) = +\infty\) and \(\varphi_{\mathrm{A}}(x) = 0\) ).

COROLLARY 4. --- For every sequence \((f_{n})\) of measurable functions \(\geqslant 0\) on X,

\[
\int^ {*} \left(\sum_ {n} f _ {n}\right) d | \mu | = \sum_ {n} \int^ {*} f _ {n} d | \mu |.\tag{1}
\]

In view of §1, No.~3, Th. 3, we are reduced to proving that for any two measurable functions \(f \geqslant 0\) , \(g \geqslant 0\) on \(X\) ,

\[
\int^ {*} (f + g) d | \mu | = \int^ {*} f d | \mu | + \int^ {*} g d | \mu |.\tag{2}
\]

This is nothing more than the additivity of the integral when \(f\) and \(g\) are integrable. In the contrary case, if for example \(f\) is not integrable, then \(\int^{*} f d|\mu| = +\infty\) by Th. 5; a fortiori, \(\int^{*}(f + g) d|\mu| = +\infty\) .

COROLLARY 5. --- For every sequence \((\mathbf{A}_{n})\) of pairwise disjoint measurable sets,

\[
| \mu | ^ {*} \left(\bigcup_ {n} \mathrm{A} _ {n}\right) = \sum_ {n} | \mu | ^ {*} (\mathrm{A} _ {n}).
\]

This follows from Cor. 4 applied to the \(\varphi_{\mathrm{A}_n}\) .
% semantic-fragments: legacy-input-seam
\relax{}
